\section{Summary}
Reproducible analysis requires a clear link between experimental design,
quality control, statistical modeling, and reporting. Here, we present a worked
Cell Reports-style analysis using a synthetic high-content screen that contains
known treatment, dose, time, and batch effects. The workflow records exclusions
before inference, normalizes measurements within plates, estimates treatment
effects with a hierarchical model, controls the false-discovery rate, and tests
whether conclusions change under alternative preprocessing choices. The
synthetic ground truth allows each analytic decision to be evaluated without
making claims about biological specimens. Across 12 simulated plates, the
pipeline recovered 92\% of implanted strong effects at a 5\% false-discovery
threshold while preserving the nominal error rate for null features. Estimated
effects were concordant across two normalization strategies
($r_{\mathrm{s}}=0.94$). This worked example shows how a compact manuscript can
make provenance, uncertainty, and robustness visible from the main text through
STAR Methods. All reported values are illustrative, and the source specification
needed to regenerate them is described in the resource-availability statement.
